\section{循环神经网络（RNN）}

\subsection{RNN 的核心思想}

循环神经网络（Recurrent Neural Network, RNN）是一种专为处理序列数据设计的神经网络架构。其核心思想是：\textbf{在每个时间步应用相同的权重矩阵}，并维护一个\textbf{隐藏状态}（hidden state）来编码已处理过的所有历史信息。

\begin{figure}[H]
\centering
\includegraphics[width=0.95\textwidth]{slides-images/slide-030.jpg}
\caption{RNN 的核心思想：同一组权重 $W$ 在每个时间步重复应用。输入序列可以是任意长度\protect\footnotemark}
\end{figure}
\footnotetext{Slides 第31页。}

\subsection{RNN 的数学定义}

在每个时间步 $t$，RNN 执行以下计算：

$$h^{(t)} = \tanh(W_h \cdot h^{(t-1)} + W_e \cdot e^{(t)} + b)$$

\begin{itemize}
    \item $h^{(t)} \in \mathbb{R}^{d_h}$：时间步 $t$ 的隐藏状态
    \item $h^{(t-1)} \in \mathbb{R}^{d_h}$：前一时间步的隐藏状态
    \item $W_h \in \mathbb{R}^{d_h \times d_h}$：隐藏状态到隐藏状态的权重矩阵
    \item $e^{(t)} \in \mathbb{R}^{d_e}$：时间步 $t$ 的词嵌入
    \item $W_e \in \mathbb{R}^{d_h \times d_e}$：输入到隐藏状态的权重矩阵
    \item $b \in \mathbb{R}^{d_h}$：偏置项
    \item $\tanh$：非线性激活函数（RNN 中最常用的选择）
\end{itemize}

初始隐藏状态 $h^{(0)}$ 通常设为全零向量。

\begin{figure}[H]
\centering
\includegraphics[width=0.95\textwidth]{slides-images/slide-032.jpg}
\caption{RNN 作为语言模型的完整架构：词经过嵌入层 $E$，通过 $W_e$ 和 $W_h$ 更新隐藏状态，最后通过 $U$ 和 softmax 预测下一个词\protect\footnotemark}
\end{figure}
\footnotetext{Slides 第33页。}

要将 RNN 用作语言模型，需要在每个时间步将隐藏状态映射到词汇表上的概率分布：

$$\hat{y}^{(t)} = \text{softmax}(U \cdot h^{(t)} + b_2)$$

\begin{importantbox}{RNN 的关键特性}
\begin{enumerate}
    \item \textbf{可处理任意长度的输入}：隐藏状态可以无限累积信息，不受固定窗口限制
    \item \textbf{参数共享}：$W_h$ 和 $W_e$ 在每个时间步都相同——无论``student''出现在序列的什么位置，都会被同一套参数处理
    \item \textbf{模型大小不随上下文增长}：无论输入多长，参数数量固定。只是计算量与序列长度成正比
\end{enumerate}
\end{importantbox}

\begin{knowledgebox}{为什么用 $\tanh$ 而不是 ReLU}
传统 RNN 中通常使用 $\tanh$ 作为激活函数，而非在前馈网络中更常见的 ReLU。$\tanh$ 的输出范围在 $[-1, 1]$，这在正负两侧是平衡的，有助于隐藏状态在多次递推后保持数值稳定。如果使用 ReLU（输出非负），隐藏状态可能会不断累加正值而爆炸。
\end{knowledgebox}

\subsection{训练 RNN 语言模型}

\subsubsection{损失函数}

RNN 语言模型的训练目标是最小化每个时间步预测的交叉熵损失：

$$J^{(t)}(\theta) = -\log P(x^{(t+1)} \mid x^{(1)}, \ldots, x^{(t)}) = -\log \hat{y}^{(t)}_{x^{(t+1)}}$$

总体损失是所有时间步损失的平均：

$$J(\theta) = \frac{1}{T} \sum_{t=1}^{T} J^{(t)}(\theta)$$

\begin{figure}[H]
\centering
\includegraphics[width=0.95\textwidth]{slides-images/slide-035.jpg}
\caption{RNN 语言模型的训练过程：在每个时间步计算预测概率分布与实际下一个词之间的交叉熵损失\protect\footnotemark}
\end{figure}
\footnotetext{Slides 第36页。}

\subsubsection{Teacher Forcing}

训练时采用\textbf{Teacher Forcing}策略：在每个时间步，无论模型预测了什么词，都将\textbf{实际的下一个词}（来自训练语料）作为下一时间步的输入。

\begin{knowledgebox}{Teacher Forcing vs 自由生成}
\textbf{Teacher Forcing}：训练时每步都``纠正''输入，使用真实的下一个词。优点是训练稳定高效；缺点是模型从未练习过从自己的错误中恢复。

\textbf{自由生成}（Free Generation/Rollout）：只在推理时使用——将模型自己生成的词作为下一步的输入，持续生成直到产生终止符 $\langle/s\rangle$。这也是 ChatGPT 等系统生成回复的方式。
\end{knowledgebox}

\subsubsection{实际训练中的分段处理}

理论上可以在整个语料上运行 RNN，但实际上这不可行——需要存储的中间状态太多。实际做法是将文本切分为固定长度的片段（如 100 个词），对每个片段独立计算损失和梯度。

\begin{warningbox}{分段引入了隐式马尔可夫假设}
虽然 RNN 理论上能处理任意长度的上下文，但分段训练意味着模型实际上只能利用片段长度以内的上下文。Manning 指出，这在某种程度上又回到了马尔可夫假设——只是窗口比 N-gram 大得多（100 词 vs 3--5 词）。现代大语言模型通过使用更长的上下文窗口（数千个 token）来缓解这一限制。
\end{warningbox}

\subsection{反向传播：BPTT}

\subsubsection{重复权重的梯度}

RNN 的一个关键特点是权重矩阵 $W_h$ 在每个时间步都被使用。那么如何计算损失对 $W_h$ 的梯度？

\begin{figure}[H]
\centering
\includegraphics[width=0.95\textwidth]{slides-images/slide-040.jpg}
\caption{RNN 的反向传播：重复权重 $W_h$ 的梯度等于它在每个时间步出现时的梯度之和\protect\footnotemark}
\end{figure}
\footnotetext{Slides 第41页。}

$$\frac{\partial J^{(t)}}{\partial W_h} = \sum_{i=1}^{t} \left. \frac{\partial J^{(t)}}{\partial W_h} \right|_{(i)}$$

\begin{importantbox}{重复权重的梯度法则}
对于重复使用的权重，其梯度等于该权重在每个位置出现时的梯度之和。数学上，这可以理解为：将 $W_h$ 在每个时间步``虚拟地''视为不同的参数 $W_h^{(1)}, W_h^{(2)}, \ldots$，分别计算梯度，然后求和。这遵循多变量链式法则中的分支求和规则。
\end{importantbox}

\subsubsection{截断反向传播}

完整的反向传播通过时间（Backpropagation Through Time, BPTT）需要展开整个序列，计算代价很高。\textbf{截断 BPTT}（Truncated BPTT）是一种实用的近似：

\begin{itemize}
    \item 前向传播仍然使用完整的上下文来更新隐藏状态
    \item 但反向传播只回溯固定的步数（如 20 步），然后截断
    \item 这大大减少了计算和存储开销，实践中通常足够好
\end{itemize}

\subsection{用 RNN 生成文本}

\begin{figure}[H]
\centering
\includegraphics[width=0.95\textwidth]{slides-images/slide-043.jpg}
\caption{RNN 文本生成的 rollout 过程：从 $\langle s \rangle$ 开始，每步采样一个词并作为下一步的输入，直到生成 $\langle/s\rangle$\protect\footnotemark}
\end{figure}
\footnotetext{Slides 第44页。}

文本生成的流程：

\begin{enumerate}
    \item 输入起始符号 $\langle s \rangle$，初始隐藏状态 $h^{(0)} = \mathbf{0}$
    \item 通过 RNN 得到下一个词的概率分布
    \item 从该分布中\textbf{采样}（sampling）一个词
    \item 将采样得到的词作为下一时间步的输入
    \item 重复步骤 2--4，直到采样到终止符 $\langle/s\rangle$
\end{enumerate}

Manning 展示了几个有趣的生成示例：

\begin{figure}[H]
\centering
\includegraphics[width=0.95\textwidth]{slides-images/slide-044.jpg}
\caption{RNN 在 Obama 演讲语料上训练后生成的文本：比 N-gram 更流畅，但仍不够连贯\protect\footnotemark}
\end{figure}
\footnotetext{Slides 第45页。}

\begin{figure}[H]
\centering
\includegraphics[width=0.95\textwidth]{slides-images/slide-045.jpg}
\caption{RNN 在 Harry Potter 语料上训练后的生成结果：能捕捉到人物名称和叙事风格\protect\footnotemark}
\end{figure}
\footnotetext{Slides 第46页。}

\begin{knowledgebox}{字符级 RNN 与条件生成}
RNN 不仅可以在词级别工作，还可以在\textbf{字符级别}逐字母生成文本。Manning 展示了一个有趣的应用：用颜色的 RGB 值初始化隐藏状态，然后让字符级 RNN 生成颜色名称（如``Ghostly Pink''、``Power Gray''等），效果出奇地好。这是\textbf{条件语言模型}（Conditional Language Model）的一个简单例子——用额外信息（颜色值）来引导生成。
\end{knowledgebox}

\subsection{RNN 的优势与缺陷}

\begin{importantbox}{RNN 相比前代模型的进步}
\begin{enumerate}
    \item 可以处理\textbf{任意长度}的输入序列
    \item \textbf{参数共享}：同一个词无论出现在什么位置，都被同样的权重处理
    \item 模型大小\textbf{不随上下文长度增长}
    \item 理论上能利用\textbf{任意远}的历史信息
\end{enumerate}
\end{importantbox}

\begin{warningbox}{RNN 的两大实际问题}
\begin{enumerate}
    \item \textbf{顺序计算，无法并行}：RNN 必须按时间步依次计算隐藏状态（本质上是一个 for 循环），无法利用 GPU 的并行计算能力。这使得 RNN 的训练速度远慢于后来的 Transformer。
    \item \textbf{长距离依赖难以学习}：虽然理论上隐藏状态包含了所有历史信息，但实际上，随着序列变长，早期的信息会在隐藏状态中逐渐``衰减''，最近的词总是支配隐藏状态。这就是所谓的\textbf{梯度消失问题}，我们将在下一节详细讨论。
\end{enumerate}
\end{warningbox}

\subsection{本章小结}

\begin{itemize}
    \item RNN 的核心思想是在每个时间步应用相同的权重，并通过隐藏状态编码历史信息
    \item 训练使用 Teacher Forcing 策略，损失函数为每步预测的交叉熵平均值
    \item 重复权重的梯度等于各时间步梯度之和；实践中常用截断 BPTT 加速
    \item RNN 能处理任意长度输入且参数共享，但受限于顺序计算和长距离依赖学习困难
\end{itemize}

%% ========== Part 4: 梯度消失与梯度爆炸 ==========

